\subsection{Small positive sources in the chosen orientation}

\begin{lemma}[Proved: a fixed-point-free P5 source on at most five states]
\label{CP2:A-P5-oriented-five}
For any finite opposite monoid $N$, a left-$N$,
right-$P_5$ biset on at most five states with no global
right fixed state cannot have a collage with the rational
cohomology ring of $\CP^2$.
\end{lemma}
\begin{proof}
There is either one regular two-point core orbit, or two.
With one orbit the two row actors agree on every state,
so Proposition~\ref{CP2:A-P5-row-independent} applies with
the one-element core-orbit set. With two orbits and no
transient, the full-core wedge lemma applies. The only
remaining case has four core states and one transient.
If the transient's two row images belong to the same
core orbit, the row actors agree and the row-independent
proposition applies with a two-element core-orbit set.
If those row images belong to different orbits, the right
act is the regular right $P_5$ act. The commuting left
action is left multiplication through a homomorphism
$N\to P_5$, and the bridge projection to $BN$ is an actual
homotopy equivalence. The collage is then homotopy
equivalent to $BP_5$. None of these arguments assumes
that $BN$ is rationally acyclic.
\end{proof}

\begin{lemma}[Proved: the nilpotent six-source low-degree comparison]
\label{CP2:A-six-nilpotent-low-degree-any-opposite}
Let $M$ be the six-element source with constant outside
actors and $a^2=e_{00}$. Suppose its right act has a
unique regular two-point band core $\{c_0,c_1\}$, with
any finite commuting opposite monoid $N$. Use
$D=\operatorname{span}\{[s]-[c_0]:s\ne c_0\}$,
$d=[c_1]-[c_0]$, and $H=A-E$ as in
Theorem~\ref{CP2:A-six-one-row-four-states}.
For its collage $X$ there is an isomorphism
\[
 H_2(X;\mathbb Q)\cong(D/HD)_N,
\]
under which the image of a generator of $H_2(BM;\mathbb Q)$
is the class of $d$, up to sign. No acyclicity or
projectivity assumption on $N$ is required here.
\end{lemma}
\begin{proof}
Use the exact all-degree source resolution displayed in
Theorem~\ref{CP2:A-six-one-row-four-states}. After tensoring
with $\mathbb Q[S]$, its degree-zero term is the trivial
$N$ module $\mathbb Qc_0$, its first differential is zero,
and its positive-degree subcomplex $C^+$ is
\[
 C^+_1=D,\qquad C^+_k=\mathbb Q\oplus D\ (k\ge2),
 \qquad d_2(t,z)=td+Hz.
\]
The right-action augmentation to a point is the projection
onto the degree-zero summand. It models the actual bridge
projection to $BN$ after taking the other bar. Thus, in
the chain homotopy pushout for the collage, that summand
cancels the $BN$ chains. What remains is the cone of
\[
 \mathbb Q\otimes_{\mathbb QN}^{\mathbf L}C^+
 \longrightarrow \widetilde C_*(BM;\mathbb Q).
\]
This is a comparison of the two actual projection maps:
the map to the source corner is induced by
$\mathbb Q[S]\to\mathbb Q$ and the augmentation of the
$N$ bar, while the map to the opposite corner is the
degree-zero projection just described.

The reduced source complex from that resolution has
$\mathbb Q$ in every degree at least two and zero
differentials. Its degree-two generator receives the
first coordinate of $C^+_2$. Using the free $N$ bar in
degrees zero and one, the degree-two group in the cone
is $\mathbb Q\oplus D$, its outgoing differential is zero,
and its degree-three boundaries span exactly
\[
 (t,-td-Hz),\qquad (0,nz-z)
 \quad(t\in\mathbb Q,\ z\in D,\ n\in N).
\]
There are no other degree-three terms contributing a
boundary in this degree. Eliminating the first coordinate
identifies the quotient with $(D/HD)_N$ and sends the
source vector $(1,0)$ to $[d]$. This proves both claims
without a high-degree comparison to $BN$.
\end{proof}

\begin{proposition}[Proved: a six-element source cannot support the chosen class on four states]
\label{CP2:A-six-oriented-four}
Let $M$ have six elements, let $N$ be any finite monoid,
and let $S$ be a left-$N$, right-$M$ biset on at most four
states. There is no projective-plane collage whose
degree-two generator restricts nontrivially to $BM$.
\end{proposition}
\begin{proof}
A global right fixed state would make that restriction
zero, so assume there is none. The positive-source
classification gives the five monoids of
Lemma~\ref{CP2:A-six-positive-shapes}.
If the band core has two regular orbits, all four states
are core states. The equivariant full-core nullhomotopy
from the proof of Theorem~\ref{CP2:A-six-one-row-four-states}
gives a $BM$ wedge summand, whose one-dimensional second
cohomology has zero square. This already excludes a
projective-plane collage.

Thus the core is the unique pair $U=\{c_0,c_1\}$.
The three split-algebra cases in that theorem give actual
rational cohomology-ring equivalences to $P_5$ collages
on quotient sets of size at most four. Those quotient
sets retain the two distinct core points and have no
right fixed state. The comparisons are $N$ linear before
taking its derived coinvariants, so they work for every
$N$. Lemma~\ref{CP2:A-P5-oriented-five} excludes them.

For the simultaneous-interchange source, use its exact
resolution $Re\to Rp_-\to Re$ and the same semilinear
row-swap lift. The core right act is $eM$, a retract of
the regular right act, so its action category is
contractible. Since $N$ fixes both core states, its core
bridge projects by a homotopy equivalence to $BN$.
The full-core nullhomotopy at the source consequently
identifies the core subcollage $X_U$ with $BM$.
The relative permutation module $\mathbb Q[S]/\mathbb Q[U]$
is annihilated by the band, so tensoring the source
resolution leaves only its $p_-$ part in degree one.
The row-swap lift acts there by $-1$.
As in Proposition~\ref{CP2:A-P5-row-independent}, all
relative homology of $(X,X_U)$ therefore has sign $-1$.
The positive fourth-homology eigenspace comes from
$H_4(BM;\mathbb Q)=0$. A square from one-dimensional
second cohomology has sign $+1$ and must vanish.

It remains to consider the nilpotent constant source.
The set-map classification from the six-source theorem
is exhaustive: $H=0$, $HD=\mathbb Qd$, or the four-state
chain $Ax=w$, $Aw=Ac_0=Ac_1=c_0$.
If $H=0$, the right transformation image is exactly the
identity and the two distinct constants, excluded by
Lemma~\ref{CP2:A-two-constant-image} for arbitrary $N$.
If $HD=\mathbb Qd$, the preceding low-degree lemma makes
$H_2(BM)\to H_2(X)$ zero. Duality contradicts the
specified nonzero restriction of the second generator.

In the chain case every opposite actor is one of
$I,F,P_0,P_1$ listed in the six-source proof. If both
$P_0,P_1$ occur, their relations kill both basis classes
in $(D/HD)_N$, so $H_2(X;\mathbb Q)=0$.
If exactly one occurs, it is a central proper biset
idempotent onto $U$. The cut-retraction theorem and the
unconditional two-cross-arrow lemma exclude this case.
If neither occurs, only the identity occurs and the
opposite action is trivial. The actual collapse to the
ordinary $M$ act cone induces an isomorphism onto second
rational cohomology by
Lemma~\ref{CP2:A-trivial-biset-side-H2}. All second
products on that cone vanish by the unconditional
order-seven act theorem. Thus its pullbacks exhaust
$H^2(X)$ and have zero square. This proves the proposition.
\end{proof}

\begin{corollary}[Proved: twenty-two remaining jointly oriented profiles]
\label{CP2:A-two-object-eighteen-oriented-twenty-two}
Under the nonzero-source-restriction orientation of
Lemma~\ref{CP2:A-two-object-positive-orientation}, the
ordered eighteen-arrow profiles $(5,8,5)$, $(5,9,4)$,
and $(6,8,4)$ are analytically excluded. The remaining
finite scope consists of the other twenty-two profiles
in Proposition~\ref{CP2:A-two-object-eighteen-scope}.
All of their opposite corners have order at most seven.
All permitted opposite tables, including the four reverse
$P_5$ target pieces, remain required.
\end{corollary}
\begin{proof}
Apply the five-state $P_5$ lemma to the first two profiles,
and the oriented six-source proposition to the third.
The complete twenty-five-profile list contains no other
opposite corner of order at least eight. These exclusions
use the chosen nonzero source restriction and are joined
through that orientation lemma; they do not assert that
an arbitrarily pre-oriented source has a nonzero
restriction.
\end{proof}
